\documentclass[9pt]{beamer}
%Information to be included in the title page:
% \title{Self-Consistency Constraints on Policy Updates:\\Axiomatic Foundations and Extensions of Mirror Descent for RL}
\title{Adaptive Step Sizes for KL Mirror Descent in RL}
\author{Philip Nielsen \and Ike Uchendu \and Kevin He \and Allen Yang}
\institute{CS2824}
\date{Spring 2026}


\newcommand{\KL}{\mathrm{KL}}
\newcommand{\Var}{\mathrm{Var}}

\usepackage{tikz}
\definecolor{accent}{HTML}{1565C0}
\definecolor{dim}{HTML}{BDBDBD}
\definecolor{cAdv}{HTML}{C2185B}   % rose
\definecolor{cVis}{HTML}{2E7D32}   % green
\definecolor{cKL}{HTML}{E65100}    % orange
\definecolor{cBud}{HTML}{1565C0}   % blue
\definecolor{cAdv}{HTML}{C2185B}
\definecolor{cVis}{HTML}{2E7D32}
\definecolor{cKL}{HTML}{E65100}
\definecolor{cBud}{HTML}{1565C0}
\definecolor{cEta}{HTML}{6A1B9A}


% Current section
\setbeamertemplate{section in toc}{%
  \vspace{0.8em}%
  \tikz[baseline=(N.base)]\node[fill=accent, text=white, rounded corners=2pt,
    inner sep=4pt, minimum width=1.8em, minimum height=1.8em,
    font=\bfseries\large] (N) {\inserttocsectionnumber};%
  \hspace{0.6em}\color{black!85}\Large\inserttocsection\par
}

% Non-current sections (dimmed)
\setbeamertemplate{section in toc shaded}{%
  \vspace{0.8em}%
  \tikz[baseline=(N.base)]\node[fill=dim, text=white, rounded corners=2pt,
    inner sep=4pt, minimum width=1.8em, minimum height=1.8em,
    font=\bfseries\large] (N) {\inserttocsectionnumber};%
  \hspace{0.6em}\color{black!40}\Large\inserttocsection\par
}

% Auto-show at every section boundary
\AtBeginSection[]{
  \begin{frame}{Outline}
  \vfill
  \begin{center}\begin{minipage}{0.85\textwidth}
  \tableofcontents[currentsection]
  \end{minipage}\end{center}
  \vfill
  \end{frame}
}



\begin{document}

\frame{\titlepage}

% \begin{frame}{Outline}
% \tableofcontents
% \end{frame}

% \begin{frame}{Outline}
% \vfill
% \begin{center}
% \begin{minipage}{0.85\textwidth}
% \tableofcontents
% \end{minipage}
% \end{center}
% \vfill
% \end{frame}


\section{Background}
\begin{frame}{Background: KL Mirror Descent for RL}


The standard KL-regularized derivation:
\vspace{0.8em}

\[
\resizebox{\textwidth}{!}{$\displaystyle
\pi_{k+1} \;=\; \arg\max_{\pi}\;
\mathbb{E}_{s\sim {\color{cVis}d_{\pi_k}}}
\!\left[
\mathbb{E}_{a\sim\pi(\cdot|s)}\!\left[{\color{cAdv}A_k(s,a)}\right]
\right]
\quad\text{s.t.}\quad
\mathbb{E}_{s\sim {\color{cVis}d_{\pi_k}}}
\!\left[{\color{cKL}\KL(\pi \,\|\, \pi_k)}\right]
\;\le\; {\color{cBud}\Delta}
$}
\]

\vspace{1.5em}

\begin{flushleft}
\small
\renewcommand{\arraystretch}{1.4}
\begin{tabular}{rl}
{\color{cAdv}$A_k(s,a)$} & Advantage: how much better is action $a$ on average \\
{\color{cVis}$d_{\pi_k}$} & State visitation under current policy \\
{\color{cKL}$\KL(\pi \,\|\, \pi_k)$} & How far the new policy has drifted \\
{\color{cBud}$\Delta$} & Budget to drift from the current policy
\end{tabular}
\end{flushleft}


\end{frame}


\begin{frame}{Background: Lagrangian Formulation}
\centering
\vspace{1em}

Move the constraint into the objective with dual variable $\lambda \ge 0$:
\vspace{0.8em}

\[
\resizebox{\textwidth}{!}{$\displaystyle
\mathcal{L}(\pi, \lambda)
\;=\;
\mathbb{E}_{s\sim {\color{cVis}d_{\pi_k}}}
\!\left[
\mathbb{E}_{a\sim\pi(\cdot|s)}\!\left[{\color{cAdv}A_k(s,a)}\right]
\right]
\;-\;
\lambda\!\left(
\mathbb{E}_{s\sim {\color{cVis}d_{\pi_k}}}
\!\left[{\color{cKL}\KL(\pi \,\|\, \pi_k)}\right]
\;-\; {\color{cBud}\Delta}
\right)
$}
\]

\vspace{1.5em}

\begin{flushleft}
\small
\renewcommand{\arraystretch}{1.4}
\begin{tabular}{rl}
$\lambda$ & Dual variable: how aggressively we enforce the budget \\
{\color{cEta}$\eta_k \;=\; 1/\lambda$} & Step size (inverse temperature)
\end{tabular}
\end{flushleft}
\end{frame}



\begin{frame}{Background: Closed-Form Solution}
\vspace{1em}

\begin{itemize}
  \setlength\itemsep{0.6em}
  \item Problem decomposes per state (KL is additive across $s$)
  \item Strictly concave in $\pi(\cdot|s)$ $\;\Rightarrow\;$ unique maximizer
  \item Set $\;\partial \mathcal{L}\,/\,\partial \pi(a|s) = 0$, solve, normalize
\end{itemize}

\vspace{1.5em}
\centering
We arrive at the \emph{exponential tilt}:
\vspace{0.6em}

\[
\boxed{\;\pi_{k+1}(a \mid s) \;\propto\; \pi_k(a \mid s)\,\exp\!\left({\color{cEta}\eta_k}\,{\color{cAdv}A_k(s,a)}\right)\;}
\]
\end{frame}



% \section{Motivation}
% \begin{frame}{Motivation}
    
% \end{frame}

\section{Motivation}
\begin{frame}{Motivation}
The KL-regularized policy update has a closed form:
\[
\pi_{k+1}(a \mid s) \;\propto\; \pi_k(a \mid s)\, \exp\!\left({\color{cEta}\eta_k}\,{\color{cAdv}A_k(s,a)}\right)
\]

\vspace{1em}
The update form is fixed. The only free parameter is the step size {\color{cEta}$\eta_k$}.

\vspace{1em}
In practice, {\color{cEta}$\eta_k$} is hand-tuned:
\begin{itemize}
\item Too large $\Rightarrow$ new policy drifts beyond where $d_{\pi_k}$ and $A_k$ are reliable
\item Too small $\Rightarrow$ wasted steps, slow convergence
  \item Good values are problem-dependent
\end{itemize}

\vspace{1em}
\textbf{Question:} are there principled ways of choosing {\color{cEta}$\eta_k$} automatically?
\end{frame}


% \begin{frame}{Motivation}
    
% \end{frame}
\begin{frame}{Motivation: Two Principled Approaches}
\vspace{1em}

We present two alternative ways to set {\color{cEta}$\eta_k$} automatically:

\vspace{2em}

\begin{center}
\begin{minipage}{0.85\textwidth}
\raggedright

\textbf{\large 1. Adaptive {\color{cEta}$\eta_k$} via KL Budget} \\[0.3em]
\hspace{1.5em}Fix how far the policy is allowed to drift each update.

\vspace{1.5em}

\textbf{\large 2. Adaptive {\color{cEta}$\eta_k$} via Distribution Shift} \\[0.3em]
\hspace{1.5em}Pick {\color{cEta}$\eta_k$} to maximize actual policy improvement.

\end{minipage}
\end{center}
\end{frame}



\section{Adaptive step size via KL budget}
\begin{frame}{Adaptive step size via KL budget}
\small
\textbf{Idea.}
Rather than fixing $\eta$ directly, choose it so that each update uses a prescribed \emph{KL budget} $\Delta$:
\[
\mathbb{E}_{s\sim d_{\pi_k}}\!\left[\KL(\pi_{\eta_k}\,\|\,\pi_k)\right]\approx \Delta.
\]
The KL divergence has exact form
\[
\KL(\pi_\eta\|\pi)=\eta\,\mathbb{E}_{\pi_\eta}[A]-\log \mathbb{E}_\pi[e^{\eta A}],
\]
and expanding the cumulant generating function to second order, we get
\[
\KL(\pi_\eta\|\pi)=\frac{\eta^2}{2}\Var_\pi(A)+O(\eta^3).
\]

\medskip
\textbf{Result.}
Using the quadratic approximation gives the variance-normalized rule
\[
\boxed{
\eta_k
=
\sqrt{
\frac{2\Delta}{
\mathbb{E}_{s\sim d_{\pi_k}}
[\Var_{\pi_k}(A_k)]
}}
}
\]
\end{frame}

\begin{frame}{Adaptive step size via KL budget (cont.)}
\small

\textbf{Why this is useful.}
Substituting the adaptive rule back into the tilt gives
\[
\pi_{k+1}(a\mid s)\propto
\pi_k(a\mid s)
\exp\!\left(
A_k(s,a)
\sqrt{
\frac{2\Delta}{
\mathbb{E}_{s\sim d_{\pi_k}}[\Var_{\pi_k}(A_k)]
}}
\right).
\]
This has several practical advantages:
\begin{itemize}
    \item \textbf{automatic normalization:} the update adapts to the current scale of the advantage signal;
    \item \textbf{fewer brittle choices:} we tune a KL budget $\Delta$ rather than a raw step size whose effect changes over training.
\end{itemize}

\medskip
However, the KL-budget argument is not yet an optimality calculation:
\begin{itemize}
    \item it does not allow for larger KL step sizes when tilting signal is clear;
    \item it treats the update locally in policy space and does \emph{not} model how changing the policy alters future state visitation and hence the downstream advantage landscape.
\end{itemize}
\end{frame}



\section{Adaptive Step Size via Distribution Shift}
\begin{frame}{Distribution Shift: Setup}
\small
\textbf{Idea.} Choose ${\color{cEta}\eta}$ to maximize the actual expected policy improvement, accounting for the fact that the tilt changes state visitation.

\medskip
Expand the value under the tilted policy to second order in $\eta$:
\[
V^{\pi_\eta}(s) = V^\pi(s)
+ \eta\,\mathrm{Cov}_{a\sim\pi(s)}(Q^{\pi_\eta},A^\pi)
+ \tfrac{\eta^2}{2}\,\mathrm{Cov}_{a\sim\pi(s)}(Q^{\pi_\eta},(A^\pi)^2)
+ O(\eta^3).
\]

\medskip
Define the improvement $\Delta V(s;\eta) := V^{\pi_\eta}(s) - V^\pi(s)$.

\medskip
\textbf{Problem.} $Q^{\pi_\eta}$ is unknown: the tilt changes state visitation, so $Q$-values under $\pi$ become stale.

\medskip
\textbf{Bellman recursion for $Q^{\pi_\eta}$.}
\begin{align*}
Q^{\pi_\eta}(s,a)
&= \mathbb{E}_{s'|s,a}[r(s,a,s') + \gamma V^{\pi_\eta}(s')] \\
&= \mathbb{E}_{s'|s,a}[r(s,a,s') + \gamma(V^\pi(s') + \Delta V(s';\eta))] \\
&= Q^\pi(s,a) + \gamma\,\mathbb{E}_{s'|s,a}[\Delta V(s';\eta)].
\end{align*}
\end{frame}


\begin{frame}{Distribution Shift: Substituting into the Covariances}
\scriptsize
Using $\mathrm{Cov}_{a}(Q^\pi,A^\pi) = \mathbb{E}_a[(A^\pi)^2]$ (since $V^\pi$ is state-only and $\mathbb{E}_a[A^\pi]=0$):
\begin{align*}
\mathrm{Cov}_{a}(Q^{\pi_\eta},A^\pi)
&= \mathrm{Cov}_{a}(Q^\pi,A^\pi) + \gamma\,\mathrm{Cov}_{a}\!\left(\mathbb{E}_{s'|s,a}[\Delta V(s';\eta)],\, A^\pi\right) \\
&= \mathbb{E}_a[(A^\pi)^2] + \gamma\,\mathrm{Cov}_{a}\!\left(\mathbb{E}_{s'|s,a}[\Delta V(s';\eta)],\, A^\pi\right).
\end{align*}
The leading term of $\Delta V$ is $\eta\,\mathbb{E}_{a'}[(A^\pi)^2] + O(\eta^2)$, so
\[
\mathrm{Cov}_{a}(Q^{\pi_\eta},A^\pi)
= \mathbb{E}_a[(A^\pi)^2]
+ \eta\gamma\,\mathrm{Cov}_{a}\!\left(\mathbb{E}_{s'|s,a}[\mathbb{E}_{a'}[(A^\pi)^2]],\, A^\pi\right)
+ O(\eta^2).
\]

\medskip
For the quadratic coefficient only the leading piece survives:
\[
\mathrm{Cov}_{a}(Q^{\pi_\eta},(A^\pi)^2)
= \mathrm{Cov}_{a}(Q^\pi,(A^\pi)^2) + \gamma\,\mathrm{Cov}_{a}\!\left(\mathbb{E}_{s'|s,a}[\Delta V(s';\eta)],\, (A^\pi)^2\right)
= \mathbb{E}_a[(A^\pi)^3] + O(\eta).
\]

\medskip
Substituting back,
\[
\Delta V(s;\eta)
= \eta\,\mathbb{E}_a[(A^\pi)^2]
+ \eta^2\gamma\,\mathrm{Cov}_{a}\!\left(\mathbb{E}_{s'|s,a}[\mathbb{E}_{a'}[(A^\pi)^2]],\, A^\pi\right)
+ \tfrac{\eta^2}{2}\,\mathbb{E}_a[(A^\pi)^3]
+ O(\eta^3).
\]
\end{frame}


\begin{frame}{Distribution Shift: Global Statistics and the Optimum}
\small
$\eta$ cannot depend on $s$, so take expectation over the state visitation $d_\pi$. Three global quantities emerge:
\begin{align*}
v^\pi        &:= \mathbb{E}_{s\sim d_\pi}\!\left[\mathbb{E}_a[(A^\pi)^2]\right]
              && \text{advantage variance} \\
c^\pi        &:= \mathbb{E}_{s\sim d_\pi}\!\left[\mathrm{Cov}_{a}\!\left(\mathbb{E}_{s'|s,a}[\mathbb{E}_{a'}[(A^\pi)^2]],\, A^\pi\right)\right]
              && \text{distribution-shift covariance} \\
\kappa_3^\pi &:= \mathbb{E}_{s\sim d_\pi}\!\left[\mathbb{E}_a[(A^\pi)^3]\right]
              && \text{advantage skewness}
\end{align*}

\medskip
Expected improvement is quadratic in $\eta$:
\[
\mathbb{E}_{s\sim d_\pi}[\Delta V(s;\eta)]
\;\approx\;
\eta\,v^\pi + \eta^2\!\left(\gamma\,c^\pi + \tfrac{\kappa_3^\pi}{2}\right).
\]

\medskip
Maximizing,
\[
\boxed{\;{\color{cEta}\eta^*} \;=\; -\dfrac{v^\pi}{2\gamma\,c^\pi + \kappa_3^\pi}\;}
\]

\medskip
\textbf{Reading.} Numerator = benefit of tilting (how informative the signal is). Denominator = penalty from one-step distribution shift ($c^\pi$) plus the nonlinearity of the tilt ($\kappa_3^\pi$).
\end{frame}


\begin{frame}{Practical Estimation: Dimensionless Normalization}
    \emph{The Scaling Problem}
    \begin{itemize}
        \item Under reward rescaling ($r \mapsto \lambda r$), variables scale at different rates: variance $v^\pi \mapsto \lambda^2 v^\pi$, but covariance $c^\pi$ and skewness $\kappa_3^\pi \mapsto \lambda^3 (\cdot)$.
    \end{itemize}

    \vspace{0.5em}
    \emph{The Dimensionless Solution}
    \begin{itemize}
        \item To stabilize the math, define a dimensionless shape statistic $s^\pi$
        \begin{equation*}
            s^\pi := \frac{2\gamma c^\pi + \kappa_3^\pi}{{v^\pi}^{3/2}}
        \end{equation*}
        \item This allows the optimal step size to be rewritten as:
        \begin{equation*}
            \eta^* = -\frac{1}{\sqrt{v^\pi}\,s^\pi}
        \end{equation*}
    \end{itemize}

    \vspace{0.5em}
    \emph{Batch Estimation}
    \begin{itemize}
        \item In practice, calculate the raw estimators ($\hat{v}_t, \hat{c}_t, \hat{\kappa}_{3,t}$) directly from a minibatch of size $B$ to construct the batch shape $\hat{s}_t$.
    \end{itemize}
\end{frame}
\begin{frame}{Practical Estimation: Moving Averages \& Stability}
    \emph{Two-Track Exponential Moving Averages }
    \begin{itemize}
        \item Rather than tracking all three global variables, maintain just two EMAs for the Scale ($V_t$) and the Shape ($S_t$):
        \begin{align*}
            V_t &= \beta_V V_{t-1} + (1-\beta_V)\hat{v}_t \\
            S_t &= \beta_S S_{t-1} + (1-\beta_S)\hat{s}_t
        \end{align*}
    \end{itemize}

    \vspace{0.2em}
    \emph{Why Two Tracks?}
    \begin{itemize}
        \item It allows $V_t$ (stable variance) to be tracked with a fast moving average, while $S_t$ (highly chaotic 3rd-order noise) is smoothed aggressively.
        \item The running step size is then reconstructed: $\eta_t = -1 / (\sqrt{V_t}\,S_t)$.
    \end{itemize}

    \vspace{0.5em}
    \emph{Safety Guardrail}
    \begin{itemize}
        \item To avoid catastrophic instability, only use this quadratic step when $S_t$ is safely negative. Otherwise, fall back to the variance-normalized rule.
    \end{itemize}
\end{frame}
\end{document}